\documentclass[10pt]{article}
\usepackage[a4paper,margin=2.6cm]{geometry}
\usepackage{array}
\usepackage{booktabs}
\usepackage{amsmath}
\usepackage{amssymb}
\usepackage[hidelinks]{hyperref}
\usepackage{xcolor}
\definecolor{ink2}{HTML}{6F6E66}
\setlength{\parskip}{2pt}
\newcommand{\code}[1]{\texttt{#1}}

\title{Family 7: a global, three-resolution input tensor\\
\large for a forward model of the Earth's surface state}
\author{blauewelt/earth --- design note}
\date{7 September 2026}

\begin{document}
\maketitle

\begin{abstract}
Family 7 is the first input tensor in this programme to cover the whole
planet. It records the ocean, atmosphere, land and ice surface on a
$0.25^\circ$ latitude--longitude grid from pole to pole, in fixed five-day
bins from 1982 to the end of 2024, in 54 channels drawn from six public
sources. Its distinguishing design choice is that channels are stored at
their native resolution in three separate groups rather than in one dense
array, and are read at a common $0.25^\circ$ position by a nearest-cell
lookup: a $1.9^\circ$ reanalysis field stored at $0.25^\circ$ would be sixty
copies of every number. The tensor holds $1.38\times10^{10}$ observed values
against the $2.5\times10^{9}$ of its North Atlantic predecessor, occupies
52.9\,GB, and is published with a per-file hash. On the sub-block it shares
with that predecessor the two agree to float16 precision on 99.2--99.5\% of
cells. We state the grid, the channels and their processing rules, the
missing-value semantics, the validation, and the decisions that shaped it,
so that a reader can load and interpret the files without further
documentation.
\end{abstract}

\section{Purpose}

The programme builds a forward model of the Earth's surface state: encode
what is observed at each grid cell into a compact vector, predict that vector
forward in time, and read out quantities of interest --- Atlantic
overturning transport, the El Ni\~no index, sea ice --- from the predicted
state. Its earlier tensors (families 3--5) covered a North Atlantic window,
$0$--$70^\circ$\,N and $100^\circ$\,W--$20^\circ$\,E, at monthly, five-day
and daily cadence. Two arguments drove the move to a global tensor. First,
the model size is set by the data rather than by compute at this scale: on
the North Atlantic the rolled forecast skill was flat from 7.6\,M to
206.66\,M parameters, and the remedy is more independent samples, which new
regimes supply and a longer record cannot. Second, a window edge is a
boundary the model must learn as if it were physics; on a sphere there is
none.

\section{Grid and time axis}

The horizontal grid is point-aligned at $0.25^\circ$: latitudes
$-90 + 0.25\,i$ for $i = 0,\dots,720$ (south first, 721 rows) and longitudes
$-180 + 0.25\,j$ for $j = 0,\dots,1439$ (1,440 columns), giving 1,038,240
cells. Longitude wraps: column 1439 ($179.75^\circ$\,E) is adjacent to column
0. A coarse $1^\circ$ grid of $181\times360$ points is defined the same way.
Every fourth $0.25^\circ$ point is a $1^\circ$ point, and a coarse channel is
read at a fine position $(y,x)$ by
\begin{equation}
y_1 = \min\!\big(\lfloor y/4 + \tfrac12\rfloor,\,180\big),\qquad
x_1 = \lfloor x/4 + \tfrac12\rfloor \bmod 360 ,
\end{equation}
so the two fine points on either side of a coarse point round to it and the
last two columns wrap to $-180^\circ$.

Time is a fixed five-day bin, a \emph{pentad}, counted from the epoch
1982-01-01: bin $b$ covers days $[5b,\,5b+5)$ after the epoch. There are
3,142 bins, the last opening on 2024-12-31. Bins are calendar-free --- they
do not align with months or years, and a year holds 73 of them --- which is
deliberate: a bin is an averaging window, not a date, and a model should not
be handed a calendar it can memorise.

\section{Channels}

Channels are stored in three groups, each an array of shape
$[T, H, W, C]$ in little-endian float16, C order, with a 128-byte NumPy
header. One time step of one group is therefore a contiguous byte range and
can be fetched on its own; the programme's web front end paints a global
frame from a single HTTP range read of 14.5\,MB.

\begin{table}[h]
\centering\small
\begin{tabular}{@{}llrrl@{}}
\toprule
group & grid & shape & size & content \\
\midrule
\code{g025} & $0.25^\circ$ & $[3142, 721, 1440, 7]$ & 45.67\,GB & ocean surface \\
\code{g100} & $1^\circ$ & $[3142, 181, 360, 15]$ & 6.14\,GB & atmosphere and land, every surface \\
\code{rg100} & $1^\circ$, monthly & $[252, 181, 360, 32]$ & 1.05\,GB & ocean interior, Argo \\
\bottomrule
\end{tabular}
\caption{The three channel groups. \code{rg100} has one row per month
(2004-01 to 2024-12); each is assigned to the pentad containing that month's
15th, and the list of live bins ships with the tensor.}
\end{table}

\subsection{Ocean surface, \code{g025}}

Five channels come from the GLORYS12 ocean reanalysis \cite{glorys}: current
speed, eastward and northward current, sea-surface height and the
$\log_{10}$ of mixed-layer depth, each a pentad mean of daily fields with at
least three days present. Speed and the two components are written in one
pass so that $\sqrt{u^2+v^2}$ equals the stored speed exactly. GLORYS covers
$80^\circ$\,S northward and begins in 1993, so rows 0--39 and bins 0--802 of
these channels are absent. Two channels are observations rather than
reanalysis: sea-surface temperature from OISST v2.1 \cite{oisst}, bilinearly
interpolated to the point grid and averaged over the pentad, present from
1982; and sea-ice concentration from the same product's \code{icec} field.
The sea-surface temperature channel is deliberately \emph{not} filled over
land or ice from any model (\S6). Table~\ref{tab:g025} lists the group.

\begin{table}[h]
\centering\setlength{\tabcolsep}{4pt}\footnotesize
\begin{tabular}{@{}rl>{\raggedright\arraybackslash}p{1.6cm}>{\raggedright\arraybackslash}p{2.6cm}>{\raggedright\arraybackslash}p{5.0cm}rr@{}}
\toprule
\# & channel & unit & source & processing rule & $\mu$ & $\sigma$ \\
\midrule
0 & \code{cur\_speed} & m\,s$^{-1}$ & GLORYS12 \code{uo},\code{vo} & pentad mean of daily $\sqrt{u^2+v^2}$; $\geq3$ days & 0.1580 & 0.1594 \\
1 & \code{log\_mld} & $\log_{10}$\,m & GLORYS12 \code{mlotst} & $\log_{10}$ of mixed-layer depth where $>0$ & 1.4785 & 0.3587 \\
2 & \code{ssh} & m & GLORYS12 \code{zos} & pentad mean of sea-surface height above geoid & $-$0.1668 & 0.7129 \\
3 & \code{cur\_u} & m\,s$^{-1}$ & GLORYS12 \code{uo} & pentad mean, eastward; same pass as speed & $-$0.0007 & 0.1793 \\
4 & \code{cur\_v} & m\,s$^{-1}$ & GLORYS12 \code{vo} & pentad mean, northward & 0.0059 & 0.1349 \\
5 & \code{sst} & $^\circ$C & OISST v2.1 \code{sst} & bilinear to point grid, pentad mean, $\geq3$ days; observed, never filled & 13.6150 & 11.6230 \\
6 & \code{sea\_ice} & fraction & OISST v2.1 \code{icec} & bilinear, pentad mean; valid range 0--1 trusted over the file's ``percent'' & 0.8271 & 0.2392 \\
\bottomrule
\end{tabular}
\caption{Group \code{g025}: seven ocean-surface channels at $0.25^\circ$, shape $[3142,721,1440,7]$. $\mu$ and $\sigma$ are the standardisation constants, in the channel's unit, computed over every finite value in every bin. GLORYS channels are absent south of $80^\circ$\,S and before 1993; sea ice is present only where concentration $\geq0.15$.}
\label{tab:g025}
\end{table}

\subsection{Atmosphere and land, \code{g100}}

Fifteen channels come from NCEP/NCAR Reanalysis~1 \cite{ncep} on its T62
Gaussian grid ($\approx1.9^\circ$), bilinearly interpolated to the $1^\circ$
point grid and averaged over the pentad: wind stress on the surface
(eastward and northward, sign-flipped from the file's momentum flux) and the
within-pentad standard deviation of each; 2\,m air temperature; 10\,m wind
components; surface pressure; precipitation rate and snow water equivalent,
each stored as $\log(1+x)$; volumetric soil moisture and soil temperature in
the top 10\,cm, masked to land \emph{before} regridding so the coastal
interpolation renormalises correctly; latent and sensible heat flux with the
reanalysis sign convention (positive upward); and skin temperature, defined
over every surface. Table~\ref{tab:g100} lists the group.


\begin{table}[h]
\centering\setlength{\tabcolsep}{4pt}\footnotesize
\begin{tabular}{@{}rl>{\raggedright\arraybackslash}p{2.2cm}>{\raggedright\arraybackslash}p{2.4cm}>{\raggedright\arraybackslash}p{5.0cm}rr@{}}
\toprule
\# & channel & unit & NCEP R1 file & processing rule & $\mu$ & $\sigma$ \\
\midrule
0 & \code{tau\_x} & N\,m$^{-2}$ & \code{uflx.sfc} & sign-flipped: stress on the surface, eastward & 0.0103 & 0.1246 \\
1 & \code{tau\_y} & N\,m$^{-2}$ & \code{vflx.sfc} & sign-flipped, northward & 0.0042 & 0.0951 \\
2 & \code{tau\_x\_std} & N\,m$^{-2}$ & \code{uflx.sfc} & within-pentad population $\sigma$ of daily $\tau_x$ & 0.0825 & 0.0724 \\
3 & \code{tau\_y\_std} & N\,m$^{-2}$ & \code{vflx.sfc} & same, northward & 0.0808 & 0.0721 \\
4 & \code{t2m} & $^\circ$C & \code{air.2m} & K $\to$ $^\circ$C & 4.8070 & 21.9986 \\
5 & \code{u10} & m\,s$^{-1}$ & \code{uwnd.10m} & eastward wind at 10\,m & 0.0027 & 4.3680 \\
6 & \code{v10} & m\,s$^{-1}$ & \code{vwnd.10m} & northward wind at 10\,m & 0.1557 & 3.1669 \\
7 & \code{sp} & hPa & \code{pres.sfc} & Pa $\to$ hPa & 965.0467 & 94.7640 \\
8 & \code{log\_prate} & $\log(1{+}x)$, mm\,d$^{-1}$ & \code{prate.sfc} & kg\,m$^{-2}$\,s$^{-1}\times86400$, then $\log(1+x)$ & 0.8662 & 0.7683 \\
9 & \code{log\_swe} & $\log(1{+}x)$, mm & \code{weasd.sfc} & snow water equivalent, $\log(1+x)$ & 2.0701 & 3.5659 \\
10 & \code{soilw} & fraction & \code{soilw.0-10cm} & volumetric; land-masked before regridding & 0.3040 & 0.0827 \\
11 & \code{tsoil} & $^\circ$C & \code{tmp.0-10cm} & K $\to$ $^\circ$C; land-masked before regridding & $-$9.6838 & 33.6608 \\
12 & \code{lhtfl} & W\,m$^{-2}$ & \code{lhtfl.sfc} & latent heat flux, positive upward & 62.3015 & 62.6118 \\
13 & \code{shtfl} & W\,m$^{-2}$ & \code{shtfl.sfc} & sensible heat flux, positive upward & 5.7186 & 40.1517 \\
14 & \code{skt} & $^\circ$C & \code{skt.sfc} & K $\to$ $^\circ$C; every surface, no mask, no splice & 5.0877 & 22.7194 \\
\bottomrule
\end{tabular}
\caption{Group \code{g100}: fifteen atmosphere-and-land channels at $1^\circ$, shape $[3142,181,360,15]$, all from NCEP/NCAR Reanalysis~1 daily Gaussian-grid files (T62), bilinearly interpolated to the $1^\circ$ point grid and averaged over the pentad with at least three days present, from 1982. Soil channels are absent over sea; every other channel is defined over every surface. The large $\sigma$ of \code{tsoil} and the low $\mu$ reflect that the $1^\circ$ land mask includes Antarctica and Greenland.}
\label{tab:g100}
\end{table}

\subsection{Ocean interior, \code{rg100}}

Thirty-two channels are the Roemmich--Gilson Argo product \cite{rg}:
temperature and salinity at sixteen pressure levels from 10 to 1900\,dbar,
bilinearly interpolated from the product's $1^\circ$ cell-centred grid and
set to missing outside its latitude band (about $64.5^\circ$\,S to
$79.5^\circ$\,N) rather than extrapolated toward the poles. Table~\ref{tab:rg100}
lists the levels.


\begin{table}[h]
\centering\setlength{\tabcolsep}{5pt}\footnotesize
\begin{tabular}{@{}rrrrrr@{}}
\toprule
level (dbar) & \code{rg\_t} $\mu$ ($^\circ$C) & $\sigma$ & \code{rg\_s} $\mu$ (PSU) & $\sigma$ & channel \# (t, s) \\
\midrule
10 & 17.867 & 9.804 & 34.818 & 1.058 & 0, 16 \\
30 & 17.519 & 9.679 & 34.864 & 1.032 & 1, 17 \\
50 & 16.797 & 9.408 & 34.919 & 1.014 & 2, 18 \\
100 & 14.731 & 8.363 & 35.004 & 0.961 & 3, 19 \\
150 & 13.032 & 7.113 & 35.013 & 0.867 & 4, 20 \\
200 & 11.704 & 6.022 & 34.962 & 0.777 & 5, 21 \\
300 & 9.825 & 4.693 & 34.838 & 0.661 & 6, 22 \\
400 & 8.436 & 3.884 & 34.739 & 0.595 & 7, 23 \\
500 & 7.340 & 3.339 & 34.668 & 0.551 & 8, 24 \\
700 & 5.792 & 2.609 & 34.604 & 0.488 & 9, 25 \\
900 & 4.704 & 2.082 & 34.603 & 0.452 & 10, 26 \\
1100 & 4.011 & 1.694 & 34.628 & 0.433 & 11, 27 \\
1300 & 3.408 & 1.235 & 34.639 & 0.219 & 12, 28 \\
1500 & 3.005 & 1.030 & 34.672 & 0.179 & 13, 29 \\
1700 & 2.689 & 0.888 & 34.697 & 0.149 & 14, 30 \\
1900 & 2.445 & 0.806 & 34.714 & 0.130 & 15, 31 \\
\bottomrule
\end{tabular}
\caption{Group \code{rg100}: the Roemmich--Gilson Argo product, temperature (\code{rg\_t}, channels 0--15) and salinity (\code{rg\_s}, channels 16--31) at sixteen pressure levels, shape $[252,181,360,32]$. One row per month from 2004-01 to 2024-12, each assigned to the pentad containing that month's 15th; bilinear from the product's $1^\circ$ cell-centred grid to the $1^\circ$ point grid; absent outside roughly $64.5^\circ$\,S--$79.5^\circ$\,N rather than extrapolated. Salinity $\sigma$ falls from 1.06 at 10\,dbar to 0.13 at 1900\,dbar, the widest scale range of any channel group.}
\label{tab:rg100}
\end{table}

\subsection{Normalisation}

Every channel is standardised before the float16 write, $x' = (x-\mu)/\sigma$
with $(\mu,\sigma)$ computed over all finite values in all bins, and the pairs
are stored per group. A stored value is therefore in standard deviations, not
in the channel's unit; float16 resolves a standardised value near 1 to about
$10^{-3}\,\sigma$, which is $0.01^\circ$C for sea-surface temperature and
$2\times10^{-4}$\,m\,s$^{-1}$ for currents. Anomalies are not stored: the
training code derives them at load time from a training-years-only
climatology.

\section{Missing values}

\code{NaN} means ``not observed here at this time'' and is information
rather than damage. The pattern is systematic. Ocean channels are absent over
land; GLORYS channels are absent south of $80^\circ$\,S and before 1993; soil
channels are absent over sea; the Argo group is absent outside 2004--2024,
outside its latitude band, and in any bin not on its monthly list. One case
is easy to misread: \textbf{sea-ice concentration is absent wherever it is
below 15\%}, because that is how the source reports ice. Measured on the
published bytes for the bin opening 2015-01-03, 155,141 of 1,038,240 cells
are finite, every one in $[0.150, 1.000]$, none zero; a mid-Atlantic cell
reads a sea-surface temperature of $21.3^\circ$C and no ice value. A
consumer wanting zeros on open water sets the channel to zero where
sea-surface temperature is finite and ice is not. Overall, $1.069\times10^{10}$
of the ocean group's $2.28\times10^{10}$ slots are finite; the three groups
together hold $1.38\times10^{10}$ observed values.

\section{Statics and truth series}

Two static fields ship with the tensor. A surface class per $0.25^\circ$
cell --- 0 ocean (702,642 cells), 1 land (226,495), 2 ice sheet or glacier
(107,074, from Natural Earth's glaciated-area polygons \cite{ne}), 3 inland
water (2,029, from its lake polygons) --- with ice taking precedence over
ocean where both apply, because the surface an instrument sees on an ice
shelf is ice. And elevation in metres from ETOPO~2022 \cite{etopo}, averaged
over the cell, negative under the sea. Two truth series are carried as
(bin, value) pairs: the RAPID array's overturning transport at
$26.5^\circ$\,N in Sverdrups, 1,459 pentads from 2004 to 2024, and the Florida
Current cable transport, 2,490 pentads from 1982. Both are Atlantic; they
exist so that a global model is scored on the same instrument as its
predecessors.

\section{Design decisions}

\paragraph{Native resolution, three groups.} The alternative, one dense
$[3142, 721, 1440, C]$ array, would store each $1.9^\circ$ reanalysis value
about sixty times and reach 425\,GB at the channel count the next phase
needs. The cost of the chosen layout is a lookup at read time; the
programme's samplers already read coarse channels at the nearest coarse cell,
so the lag-0 $3\times3$ patch of a fine position serves the same coarse value
nine times, which is the model's honest view of a coarse field rather than a
rounding artefact.

\paragraph{Two temperatures, split by instrument.} An earlier build merged
observed sea-surface temperature with reanalysis skin temperature into one
channel. The merge was reversed: it spliced an infrared--microwave analysis
over the sea onto a reanalysis everywhere else at the coastline, which the
model would have had to learn as an instrument boundary. The rule adopted is
that a channel is shared across surfaces only when both the measurand and the
instrument match. Sea-surface temperature therefore stays observed and
absent over land; skin temperature is reanalysis everywhere. The two are
different measurements and are stored as such.

\paragraph{Sea ice and snow stay separate.} A fraction and a water mass
cannot be merged into a ``frozen fraction'' without an invented constant.

\paragraph{The atmosphere source is a stand-in.} NCEP/NCAR Reanalysis~1 was
chosen over ERA5 because it needs no account; ERA5 replaces the
\code{g100} fetch in a later phase with no change of layout.

\paragraph{Resumable, hash-verified publication.} The build runs as one job
resumable per stage, with a per-stage specification digest so that a changed
stage discards only its own state, and publishes to a public dataset
repository with a per-file SHA-256 that is verified by downloading each file
back. Three runs were needed --- the first cancelled for the temperature
merge above, the second failed on a path error in the truth stage --- at a
total cost of about \$1.9.

\section{Validation}

Ten assertions run before the tensor is trusted, on a CPU with no network:
row 0 is the South Pole, asserted on the bytes; $\sqrt{u^2+v^2}$ equals the
stored speed; the absence pattern south of $80^\circ$\,S; the sea-ice unit
(the source declares ``percent'' but its valid range is 0--1, and the reader
trusts the range); the surface-class counts; normalisation shapes; truth
keys; and resumability across a simulated interruption.

The tensor is not bit-identical to its North Atlantic predecessor, and was
not expected to be. On the shared sub-block at the bin opening 2015-01-03,
over the 86,698 common ocean cells, currents, sea-surface height and
mixed-layer depth agree to two float16 quanta on 99.2--99.5\% of cells
(correlation $\geq 0.99995$, median absolute difference $10^{-5}$\,m\,s$^{-1}$),
with a tail of about 1\% differing by up to 0.19\,m\,s$^{-1}$ because the
daily fields were binned in a different order at partially masked coastal
blocks. Sea-surface temperature agrees to $0.03^\circ$C. The global tensor
observes 1,756 more coastal cells than the window did, because its
regridding renormalises at the coast rather than leaving a masked neighbour
empty.

\section{Limitations}

Three limitations are stated rather than hidden. The atmosphere and land
channels are a $1.9^\circ$ reanalysis, the coarsest source in the tensor.
Subsurface ocean structure is a monthly climatological product at sixteen
levels, present only from 2004 and only within the Argo latitude band. And
the record is 43 years: no data source shortens that constraint, and it is
the reason the programme's models are small.

\section{Next steps: the data ladder}

The programme's stated aim is a model of $10^9$--$10^{10}$ parameters. At the
programme's Chinchilla anchor of one parameter per twenty observed values,
that size is earned by $2\times10^{10}$--$2\times10^{11}$ values, against the
$1.38\times10^{10}$ family 7 holds. The question is therefore which data to
add, in what order, and the answer is set by a measurement rather than by
taste: a $0.25^\circ$ cell is correlated with its neighbours over roughly
300\,km and over two to six months at the surface, so the raw value count
overstates the independent sample count by about two orders of magnitude,
and the flatness of rolled skill from 7.6\,M to 206.66\,M parameters on the
North Atlantic is the symptom. Data that adds \emph{new regimes} or
\emph{new physics} is worth its face value; data that adds the same sample
five times is not. That ordering is: other spheres first, then depth, then
cadence, and horizontal resolution last.

\begin{table}[h]
\centering
\setlength{\tabcolsep}{4pt}\footnotesize
\begin{tabular}{@{}>{\raggedright\arraybackslash}p{2.5cm}>{\raggedright\arraybackslash}p{5.9cm}>{\raggedright\arraybackslash}p{1.9cm}>{\raggedright\arraybackslash}p{1.2cm}>{\raggedright\arraybackslash}p{3.0cm}@{}}
\toprule
rung & what is added & values & cost & prerequisite \\
\midrule
7 (today) & the global surface at $0.25^\circ$, atmosphere and land at $1^\circ$ & $1.4\times10^{10}$ & 53\,GB & --- \\
\addlinespace
L1: the other spheres, properly & ERA5 replacing the NCEP channels at the same layout; observed land-surface temperature (MODIS); scatterometer winds (ASCAT); soil moisture (SMAP, CCI); terrestrial water storage (GRACE); sea-ice concentration and drift (OSI SAF) & $\approx9\times10^{10}$ & $\approx$170\,GB & a Copernicus Climate Data Store account, which only the project owner can create \\
\addlinespace
depth & GREP three-dimensional ocean reanalysis at $0.25^\circ$: four variables at eight levels, subsurface velocity and abyssal levels the tensor has never had, eleven extra years over Argo & $+20$\,\% & $+20$\,GB & existing credentials \\
\addlinespace
observed surface overturning & DUACS altimetry (absolute dynamic topography, geostrophic currents) and OSTIA foundation SST, both taken at $0.25^\circ$ & small & $+4.7$\,GB & existing credentials \\
\addlinespace
daily cadence & every group at one-day bins & $\times5$ raw, $\approx{\times}1.5$ indep. & 1.3\,TB & a streaming loader; no tensor fits a box \\
\addlinespace
the kilometre rung & 1\,km daily radiometry, fed to a patch encoder rather than stored as a tensor & $3.5\times10^{14}$ & --- & the first rung that is compute-bound \\
\bottomrule
\end{tabular}
\caption{The data ladder above family 7, in the order the programme intends
to climb it. Values are observed values; ``independent'' applies the
correlation correction of the text.}
\end{table}

\paragraph{Why the spheres come first.} The L1 rung multiplies the observed
values by about 4.5 at 170\,GB, and every value it adds is new physics: land
does not advect, the atmosphere decorrelates in days, ice has its own
seasonal cycle. Family 7 already carries the first phase of this rung --- the
fifteen NCEP channels are the L0 set, with a reanalysis standing in for ERA5
until an account exists --- so the layout, the lookup and the loader are in
place and the upgrade is a change of source, not of shape. It is also the
rung that unlocks the read-outs the programme has promised beyond the
Atlantic: a sea-ice head and a land-water head need the observed channels,
not the reanalysis ones.

\paragraph{Why depth comes before pixels.} The overturning is a vertical
circulation, and the tensor's only view below the surface is a monthly
climatological product. A daily three-dimensional reanalysis at eight levels
adds the subsurface velocity field for 20\,GB on credentials the programme
already holds, and reaches back to 1993 where Argo begins in 2004. It is
ranked first in the programme's own data-ladder document on the ratio of new
physics to new risk.

\paragraph{Why resolution comes last.} A $1/12^\circ$ basin would cost
305\,GB for the present channels and, by the correlation argument above, add
far fewer independent samples than bytes. The experiment that should precede
any such regrid is small: a single $1/12^\circ$ strip across the RAPID
section, $23$--$30^\circ$\,N, 18\,GB for forty channels, asking whether
resolution at the section changes the transport read-out at all. If it does
not, the basin regrid is never bought. Kilometre-scale radiometry is a
different object altogether --- $3.5\times10^{14}$ values cannot be a tensor
--- and enters, when it enters, as a patch-encoder input in the style of
AlphaEarth's per-source encoders.

\paragraph{Two cheap items that precede all of the above.} The 2025--2026
continuation of every source adds 4.5\% of end-bins nearly for free and lets
the terminal holdout run to 2026 instead of 2024. And the static fields the
programme downloads but does not yet use --- bathymetric slope, the Coriolis
parameter and its gradient, distance to coast --- are megabytes with the
highest information per byte of anything on this list.

\begin{thebibliography}{9}\small
\bibitem{glorys} Lellouche, J.-M.\ et al.\ (2021). The Copernicus Global
$1/12^\circ$ Oceanic and Sea Ice GLORYS12 Reanalysis. \emph{Front.\ Earth
Sci.}\ 9:698876.
\bibitem{oisst} Huang, B.\ et al.\ (2021). Improvements of the Daily Optimum
Interpolation Sea Surface Temperature (DOISST) Version 2.1. \emph{J.\
Climate} 34, 2923--2939.
\bibitem{ncep} Kalnay, E.\ et al.\ (1996). The NCEP/NCAR 40-Year Reanalysis
Project. \emph{Bull.\ Amer.\ Meteor.\ Soc.}\ 77, 437--471.
\bibitem{rg} Roemmich, D.\ and Gilson, J.\ (2009). The 2004--2008 mean and
annual cycle of temperature, salinity, and steric height in the global ocean
from the Argo Program. \emph{Prog.\ Oceanogr.}\ 82, 81--100.
\bibitem{etopo} NOAA National Centers for Environmental Information (2022).
ETOPO 2022 15 Arc-Second Global Relief Model.
\bibitem{ne} Natural Earth, version 5.1.1, 1:10m physical vectors
(glaciated areas; lakes). Public domain.
\end{thebibliography}

\end{document}
